%% kapa: exam=1 round=20 year=2009 period=1 form=A subject=civil_law subject_ko=민법 seq=5 number=5
%% source: Q-Net 감정평가사 기출문제 (cst00309) artlSeq=5024178 file=extracted/1008/5024178/2009년 제20회 감정평가사 1차시험 기출문제/1교시.hwp
%% conversion: kordoc 4.24.51 -> tools/gen1.py (auto); review=proofread-sonnet; flags=none
\begin{question}{5}
태아의 권리능력에 관한 설명으로 옳은 것은? (다툼이 있으면 판례에 의함)
\begin{choices}
  \item 태아는 증여의 상대방이 될 수 있다.
  \item 정지조건설에 따르면, 태아의 법정대리인이 필요하다.
  \item 민법규정에 따르면 태아는 부(父)에 대해 인지를 청구할 수 있다.
  \item 모체에 대한 가해행위로 모(母)와 함께 태아가 사망한 경우, 해제조건설에 따르면 태아 자신은 손해배상청구권을 가지지 못한다.
  \item 태아의 부(父)가 타인의 불법행위로 사망한 경우, 부(父)의 생명침해로 인한 재산적 손해에 대해서는 태아에게 직접 손해배상청구권이 발생한다.
\end{choices}
\end{question}
\answer{4}
